% Related Work
\section{Related Work}
\label{sec:related}

Our work intersects three research areas: AI alignment methods that optimize models to match human preferences, user adaptation studies in HCI that examine how humans learn to use AI tools, and conversation-level quality metrics that evaluate multi-turn interactions. While each area has made significant progress, no prior work measures bidirectional alignment---the coupling between AI adaptation (trained responsiveness) and user adaptation (learned query strategies) during conversations.

\subsection{AI Alignment and RLHF}

Reinforcement Learning from Human Feedback (RLHF) has become the dominant paradigm for aligning large language models with human preferences \cite{christiano2017deep, ouyang2022training}. The standard approach trains a reward model on human preference judgments (response A preferred over response B), then uses reinforcement learning to optimize the language model policy to maximize predicted reward. InstructGPT \cite{ouyang2022training} demonstrated that RLHF-trained models produce outputs humans rate as more helpful, honest, and harmless than supervised fine-tuning alone. Constitutional AI \cite{bai2022constitutional} extended this framework by incorporating AI-generated feedback to reduce reliance on human annotation.

These methods treat alignment as unidirectional optimization: adjust the AI policy until human evaluators are satisfied. Evaluation focuses on aggregate preference accuracy, single-turn response quality, or held-out test sets where humans rate AI outputs. Success is measured by how well the AI conforms to human preferences, not by how the human-AI interaction evolves. While RLHF produces models that generate high-quality individual responses, it does not account for conversation dynamics---how users learn AI capabilities over turns or how AI systems respond to evolving query patterns within extended dialogues.

Our work complements RLHF by proposing conversation-level metrics that measure alignment during interaction. Rather than asking ``does the AI produce responses humans prefer?'', we ask ``do users and AI co-adapt during multi-turn conversations?'' This shift from static policy evaluation to dynamic interaction measurement enables detection of alignment patterns that aggregate ratings may miss---for instance, conversations where users struggle to communicate effectively despite receiving high-quality individual responses.

\subsection{User Adaptation and Learning Curves}

Human-computer interaction research has long recognized that users learn to communicate more effectively with systems over time \cite{newell1981mechanisms}. Learning curve theory predicts that task performance improves with practice, typically following a power law where early interactions show rapid improvement that plateaus with experience. In the context of conversational AI, this manifests as query reformulation patterns: novice users reformulate frequently (rephrasing questions, adding clarifications) while experienced users craft effective queries from the start.

Prior work on query reformulation has focused on search engines and information retrieval systems \cite{huang2009query}. Users reformulate queries when initial results are unsatisfactory, and reformulation rates correlate with search success---effective reformulations improve retrieval quality while excessive reformulation signals user frustration. However, this literature treats the search engine as a static system; reformulation reflects user learning about the corpus and ranking algorithm, not about an adaptive AI agent.

More recent work examines how users adapt to AI assistants. Studies of voice interfaces \cite{luger2016like} find that users develop mental models of system capabilities through trial-and-error, learning which requests succeed and which fail. Conversational repair strategies \cite{ashktorab2019resilient} show users employ systematic patterns when AI misunderstands them---rephrasing, adding context, or breaking requests into sub-tasks. These findings suggest users actively learn AI capabilities during conversations, but existing work has not quantified learning rates or connected them to alignment quality.

Our contribution is to operationalize user learning through \textbf{reformulation slope}---the rate at which reformulation frequency decreases over conversation turns. Negative slope indicates learning (reformulation decreases), while flat or positive slope suggests no learning or engagement decay. By measuring slope as a conversation-level metric, we can test whether user adaptation correlates with interaction quality and whether it couples with AI responsiveness.

\subsection{Conversational AI Evaluation}

Evaluation of multi-turn conversational systems has traditionally focused on task completion, dialogue coherence, and user satisfaction \cite{walker1997paradise}. Task-oriented dialogue systems are evaluated on success rate (did the user accomplish their goal?) and efficiency (how many turns required?). Open-domain chatbots are evaluated on engagement (conversation length), appropriateness (avoiding offensive outputs), and subjective quality ratings.

Recent work has introduced conversation-level metrics beyond aggregate ratings. Mehri \& Eskenazi \cite{mehri2020usr} propose USR, an automatic metric that correlates with human judgments of naturalness and coherence. Zhang et al.\ \cite{zhang2021dialogpt} analyze turn-level topic shifts and introduce metrics for conversation depth and breadth. Durmus et al.\ \cite{durmus2023measuring} evaluate helpfulness in multi-turn contexts by measuring whether AI responses enable users to complete complex tasks.

However, these metrics still treat the AI as the primary unit of analysis---measuring response quality, coherence, or task enablement. They do not measure user behavioral changes during conversation (learning curves) or the coupling between user and AI adaptation patterns. A conversation where the AI produces individually high-quality responses but users struggle to communicate effectively (high reformulation rate, no learning) would receive high scores on existing metrics but low scores on bidirectional alignment.

Interactive Alignment Theory from psycholinguistics \cite{pickering2004toward} provides a relevant framework. In human-human dialogue, both participants unconsciously align their linguistic representations---lexical choices, syntactic structures, situational models---through a process of priming and adaptation. Successful conversations exhibit high alignment on multiple levels. While this theory describes human cognition, the core insight applies to human-AI interaction: alignment emerges through mutual adaptation during conversation, not from one party conforming to the other.

Our work extends conversation evaluation by introducing bidirectional metrics. We measure user learning (reformulation slope) alongside AI responsiveness (diversity correlation) and propose coupling strength (correlation between these signals) as a conversation-level alignment metric. This approach captures interaction dynamics that aggregate quality scores may miss, providing a complementary view of alignment focused on co-adaptation rather than static policy quality.

\subsection{Positioning Our Contribution}

To our knowledge, no prior work has operationalized bidirectional alignment through observable conversation behaviors. RLHF measures alignment as policy optimization against human preferences. HCI learning curve research measures user adaptation but treats AI as static. Conversation evaluation metrics measure output quality but not behavioral dynamics. We combine insights from all three areas: RLHF provides the training context (aligned AI systems), learning curve theory provides the user adaptation framework (reformulation slopes), and conversation analysis provides the interaction context (multi-turn dynamics). The result is a novel metric---behavioral coupling---that quantifies co-adaptation between users and AI during conversations.
